\documentclass[a4paper,12pt]{article}
\usepackage[utf8]{inputenc}
\usepackage[italian]{babel}
\usepackage{listings}
\usepackage{xcolor}
\usepackage{float}
\usepackage{placeins}

\lstset{
  language=C,
  backgroundcolor=\color{lightgray!20},
  frame=single,
  breaklines=true,
  numbers=left,
  numberstyle=\tiny,
  stepnumber=1,
  numbersep=5pt,
  captionpos=b,
  basicstyle=\ttfamily\small,
  keywordstyle=\color{blue}\bfseries,
  commentstyle=\color{green!60!black},
  stringstyle=\color{red},
  showstringspaces=false,
  float=H
}

\title{LabReSiD26 - Hands-On 3}
\author{Davide Ferrara - 518629}
\date{}

\begin{document}

\maketitle

% ---------------------------------------------------------------------------
\section{Esercizio 1: Identifica i problemi nel codice di partenza}
% ---------------------------------------------------------------------------

\subsection{1. Cosa succede se l'utente lancia \texttt{./sniffer} senza argomenti?}

La seguente linea converte il nome dell'interfaccia in indice numerico.
Lo stesso indice è visualizzabile col comando \texttt{ip a}:

\begin{lstlisting}
  /* Recupera l'indice numerico dell'interfaccia */
  unsigned int ifindex = if_nametoindex(iface);
\end{lstlisting}

Nel caso in cui ifindex fosse \textbf{null} il programma al momento della
chiamata di \texttt{if\_nametoindex(NULL)} andrebbe in Segmentation fault
in quanto proverebbe ad accedere un'area di memoria di cui non dispone i permessi.

\begin{lstlisting}[language={}]
  sudo ./sniffer
  [sudo] password for dff:
  [SNIFFER.c] Sniffing su (null)...
  Segmentation fault
\end{lstlisting}

È importante controllare anche il valore di ritorno di \texttt{if\_nametoindex()},
che ritorna 0 in caso di errore.
Se non si controlla, viene passato 0 a \texttt{bind()}, che in
\texttt{AF\_PACKET} interpreta l'indice 0 come ``tutte le interfacce'': il programma non
crasha ma sniffa silenziosamente su tutte le interfacce invece di fallire.

\begin{lstlisting}[language={}]
ip a
1: lo: <LOOPBACK,UP,LOWER_UP> mtu 65536 qdisc noqueue state UNKNOWN group default qlen 1000
    link/loopback 00:00:00:00:00:00 brd 00:00:00:00:00:00
2: enp8s0: <NO-CARRIER,BROADCAST,MULTICAST,UP> mtu 1500 qdisc fq_codel state DOWN group default qlen 1000
    link/ether 74:56:3c:c7:5e:5d brd ff:ff:ff:ff:ff:ff
4: wlan0: <BROADCAST,MULTICAST,UP,LOWER_UP> mtu 1500 qdisc noqueue state UP group default qlen 1000
    link/ether bc:c7:46:9b:2e:fe brd ff:ff:ff:ff:ff:ff
    inet 192.168.8.230/24 metric 600 brd 192.168.8.255 scope global dynamic wlan0
\end{lstlisting}

Nel mio caso l'indice di wlan0 (la scheda wifi) corrisponde a indice 4.

Senza il controllo su \texttt{ifindex}, come si puo notare snifferei su tutte le 
interfacce:

\begin{lstlisting}[language={}]
sudo ./sniffer foo
[SNIFFER.c] Sniffing su foo...
[SNIFFER.c] Indice interfaccia 0
^C
\end{lstlisting}

Con il controllo invece esco dal programma:

\begin{lstlisting}[language={}]
sudo ./sniffer foo
[SNIFFER.c] Sniffing su foo...
Errore: interfaccia 'foo' non trovata
\end{lstlisting}

\subsection{2. Perché controllare il valore di ritorno di \texttt{socket()} e \texttt{bind()}?}

Van controllate per evitare che il programma continui con un file descriptor 
non valido, causando comportamenti indefiniti alla chiamata di
(\texttt{bind()}, \texttt{recvfrom()}).

Esempi di errori di \textbf{\texttt{socket()}} che ho riscontato sono durante
la scrittura del programma sono:

\begin{itemize}
  \item \texttt{EPERM}: il processo non ha i privilegi necessari (serve \texttt{sudo} o
        la capability \texttt{cap\_net\_raw}).
\end{itemize}

\texttt{cap\_net\_raw} è una Linux capability che permette di aprire raw socket senza
essere root. Invece di eseguire l'intero programma con \texttt{sudo}, si assegna solo
questo privilegio specifico al binario:

\begin{lstlisting}[language={}, caption={Assegnare cap\_net\_raw al binario}]
sudo setcap cap_net_raw+ep ./sniffer
./sniffer wlan0   # funziona senza sudo
\end{lstlisting}

Le capability assegnate al binario sono salvate come extended attribute
(\texttt{security.capability}) sull'inode del file e leggibili con \texttt{getcap}:

\begin{lstlisting}[language={}, caption={Lettura capability con getcap}]
getcap ./sniffer
./sniffer cap_net_raw=ep
\end{lstlisting}

Gli errori invece riscontrati con bind():

\begin{itemize}
  \item \texttt{ENODEV}: l'indice di interfaccia passato non corrisponde a nessun dispositivo.
\end{itemize}

Il codice qui sotto effettua il controllo di socket() e bind() e su EPERM
stampo un messaggio personalizzato.
\begin{lstlisting}[caption={Controllo socket() e bind()}]
int sock = socket(AF_PACKET, SOCK_RAW, htons(ETH_P_ALL));
if (sock == -1) {
  if (errno == EPERM)
    fprintf(stderr, "Errore: esegui con sudo o cap_net_raw\n");
  else
    perror("socket");
  exit(EXIT_FAILURE);
}

if (bind(sock, (struct sockaddr *)&sll, sizeof(sll)) == -1) {
  perror("bind");
  close(sock);
  exit(EXIT_FAILURE);
}
\end{lstlisting}

\subsection{3. Perché \texttt{eth->h\_proto} stampato direttamente dà un valore sbagliato su x86?}

Il campo \texttt{h\_proto} arriva dalla rete in \textbf{big-endian} (byte più significativo
per primo): IPv4 è \texttt{0x0800}, cioè i byte sul cavo sono \texttt{[0x08, 0x00]}.

x86 è \textbf{little-endian}: il processore legge il byte all'indirizzo più basso come
il meno significativo. Quindi \texttt{[0x08, 0x00]} viene interpretato come
\texttt{0x0008} = 8, invece del corretto \texttt{0x0800} = 2048.

\texttt{ntohs()} (\textit{network to host short}) esegue il byte-swap necessario:

\begin{lstlisting}[caption={Senza e con conversione}]
/* SBAGLIATO: stampa 8 invece di 2048 */
printf("  EtherType: 0x%04X\n", eth->h_proto);

/* CORRETTO: ntohs() swappa i byte e stampa 2048 */
printf("  EtherType: 0x%04X\n", ntohs(eth->h_proto));
\end{lstlisting}

\subsection{4. Cosa succede se si omette la \texttt{bind()}?}

Dal manuale di \texttt{bind(2)}:

\begin{quote}
\textit{When a socket is created with socket(2), it exists in a name space
(address family) but has no address assigned to it. bind() assigns the address
specified by addr to the socket referred to by the file descriptor sockfd.
Traditionally, this operation is called "assigning a name to a socket".}
\end{quote}

Senza \texttt{bind()}, il socket \texttt{AF\_PACKET} non è associato ad alcuna interfaccia
e riceve frame da \textbf{tutte le interfacce} attive sul sistema. Con \texttt{bind()}
il kernel consegna solo i frame dell'interfaccia specificata.

Nella prova pratica la differenza non è evidente perché ricevo frame solo
su wlan0 mentre \texttt{enp8s0} è in stato \texttt{DOWN} e non genera frame. 
Su un sistema con più interfacce attive simultaneamente (es.\ \texttt{eth0}
e \texttt{wlan0}), senza \texttt{bind()} si vedrebbero frame da entrambe, mischiati e
indistinguibili.

% ---------------------------------------------------------------------------
\section{Esercizio 2: Aggiungi la gestione degli errori}
% ---------------------------------------------------------------------------

La gestione degli errori è stata già implementata durante l'analisi
dell'Esercizio 1.
Il codice risultante è riportato di seguito per completezza.

\begin{lstlisting}[caption={Gestione errori -- sniffer.c}]
/* Controlla argc */
if (argc < 2) {
  fprintf(stderr, "Uso: %s <interfaccia> [arp|ipv4|ipv6]\n", argv[0]);
  exit(EXIT_FAILURE);
}
const char *iface = argv[1];

/* Apre il raw socket */
int sock = socket(AF_PACKET, SOCK_RAW, htons(ETH_P_ALL));
if (sock == -1) {
  if (errno == EPERM)
    fprintf(stderr, "Errore: esegui con sudo o cap_net_raw\n");
  else
    perror("socket");
  exit(EXIT_FAILURE);
}

...

/* Recupera l'indice dell'interfaccia */
unsigned int ifindex = if_nametoindex(iface);
if (ifindex == 0) {
  perror("if_nametoindex");
  close(sock);
  exit(EXIT_FAILURE);
}

...

if (bind(sock, (struct sockaddr *)&sll, sizeof(sll)) == -1) {
  perror("bind");
  close(sock);
  exit(EXIT_FAILURE);
}
\end{lstlisting}

% ---------------------------------------------------------------------------
\section{Esercizio 3: Decodifica l'EtherType}
% ---------------------------------------------------------------------------

La funzione \texttt{ethertype\_name()} riceve il valore \texttt{proto} già convertito in
host byte order (dopo \texttt{ntohs()}) e restituisce la stringa corrispondente al protocollo.
Riconosce IPv4, ARP e IPv6; per qualsiasi altro valore restituisce \texttt{"sconosciuto"}.

Invece dei magic number si usano le costanti simboliche definite in
\texttt{<net/ethernet.h>}, già disponibili in host byte order:

\begin{lstlisting}[caption={Costanti EtherType da net/ethernet.h (estratto)}]
#ifndef ETHERTYPE_IP
#define ETHERTYPE_IP   0x0800  /* IP protocol */
#endif
#ifndef ETHERTYPE_ARP
#define ETHERTYPE_ARP  0x0806  /* Addr. resolution protocol */
#endif
#ifndef ETHERTYPE_IPV6
#define ETHERTYPE_IPV6 0x86dd
#endif
\end{lstlisting}

Uso semplicemente uno switch case per i vari casi:
\begin{lstlisting}[caption={ethertype\_name() -- implementazione con costanti simboliche}]
const char *ethertype_name(uint16_t proto) {
  switch (proto) {
  case ETHERTYPE_IP:
    return "IPv4";
  case ETHERTYPE_ARP:
    return "ARP";
  case ETHERTYPE_IPV6:
    return "IPv6";
  default:
    return "sconosciuto";
  }
}
\end{lstlisting}

Nel loop del programma:
\begin{lstlisting}[caption={Uso nel loop di stampa}]
uint16_t proto = ntohs(eth->h_proto);
const char *name = ethertype_name(proto);
printf("  EtherType: %d %s\n", proto, name);
\end{lstlisting}

Esempio di Output:
\begin{lstlisting}[language={}, caption={Output -- Es.3}]
Frame #1  (139 byte)
  DST: BC:C7:46:9B:2E:FE
  SRC: 94:83:C4:72:F8:D1
  EtherType: 2048 IPv4
Frame #2  (66 byte)
  DST: 94:83:C4:72:F8:D1
  SRC: BC:C7:46:9B:2E:FE
  EtherType: 2048 IPv4
\end{lstlisting}

% ---------------------------------------------------------------------------
\section{Esercizio 4: Statistiche e interruzione pulita}
% ---------------------------------------------------------------------------

Alla ricezione di \texttt{SIGINT} (Ctrl+C), l'handler imposta \texttt{running = 0}
interrompendo il loop. Dopo l'uscita dal loop si chiama \texttt{print\_stats()} che
stampa i contatori per protocollo. Il controllo \texttt{if (running == 0)} dentro il
loop permette di stampare le statistiche anche se l'interrupt arriva durante
l'\texttt{elaborazione} del frame corrente (immagino debba finire l'esecuzione
un thread? perché senno mi stampava prima le statistiche e successivamente
altre informazioni sniffate, per questo ho fatto cosi).

\begin{lstlisting}[caption={Variabili globali, handler e print\_stats()}]
volatile int running = 1;
unsigned long cnt_ipv4 = 0;
unsigned long cnt_arp   = 0;
unsigned long cnt_ipv6  = 0;
unsigned long cnt_other = 0;

void print_stats(void) {
  printf("\nSNIFFING STATS: PROTOCOL | COUNT\n");
  printf("                  IPv4     | %ld\n", cnt_ipv4);
  printf("                  ARP      | %ld\n", cnt_arp);
  printf("                  IPv6     | %ld\n", cnt_ipv6);
  printf("                  Altri    | %ld\n", cnt_other);
}

void handle_sigint(int sig) {
  (void)sig;
  running = 0;
}
\end{lstlisting}

I contatori vengono incrementati direttamente dentro \texttt{ethertype\_name()},
nei rispettivi \texttt{case}: la funzione ha quindi un side effect oltre a restituire
la stringa. Nel loop basta chiamare \texttt{ethertype\_name(proto)} e i contatori
si aggiornano automaticamente.

\begin{lstlisting}[caption={ethertype\_name() con side effect sui contatori}]
const char *ethertype_name(uint16_t proto) {
  switch (proto) {
  case ETHERTYPE_IP:   cnt_ipv4++; return "IPv4";
  case ETHERTYPE_ARP:  cnt_arp++;  return "ARP";
  case ETHERTYPE_IPV6: cnt_ipv6++; return "IPv6";
  default:             cnt_other++; return "sconociuto";
  }
}
\end{lstlisting}

\begin{lstlisting}[caption={Loop -- nessuno switch separato}]
while (running) {
  /* ... recvfrom e parsing ... */
  uint16_t proto = ntohs(eth->h_proto);
  const char *name = ethertype_name(proto); /* incrementa contatore */
  printf("  EtherType: %d %s\n", proto, name);

  if (running == 0)
    print_stats();
}
\end{lstlisting}

Esempio di Output:
\begin{lstlisting}[language={}, caption={Output -- Ctrl+C dopo 5 frame}]
sudo ./sniffer wlan0
[SNIFFER.c] Sniffing su wlan0...
[SNIFFER.c] Indice interfaccia 4
Frame #1  (139 byte)
  DST: BC:C7:46:9B:2E:FE
  SRC: 94:83:C4:72:F8:D1
  EtherType: 2048 IPv4
Frame #2  (66 byte)
  DST: 94:83:C4:72:F8:D1
  SRC: BC:C7:46:9B:2E:FE
  EtherType: 2048 IPv4
^CFrame #5  (139 byte)
  DST: BC:C7:46:9B:2E:FE
  SRC: 94:83:C4:72:F8:D1
  EtherType: 2048 IPv4

SNIFFING STATS: PROTOCOL | COUNT
                  IPv4     | 5
                  ARP      | 0
                  IPv6     | 0
                  Altri    | 0
\end{lstlisting}

% ---------------------------------------------------------------------------
\section{Esercizio 5: Filtraggio da riga di comando}
% ---------------------------------------------------------------------------

Il secondo argomento è opzionale. \texttt{filter} è una variabile \texttt{uint16\_t}
inizializzata a 0 prima del loop: con \texttt{filter = 0} il loop non salta nessun frame
(nessun filtro attivo). Se \texttt{argv[2]} è presente, viene convertito nella costante
EtherType corrispondente una sola volta prima del loop.

Nel loop basta una riga: se \texttt{filter != 0} e il protocollo del frame non corrisponde,
il frame viene scartato con \texttt{continue}.

\begin{lstlisting}[caption={Parsing del filtro da argv[2] -- prima del loop}]
uint16_t filter = 0; /* 0 = nessun filtro */
if (argc == 3) {
  if      (strcmp(argv[2], "ipv4") == 0) filter = ETHERTYPE_IP;
  else if (strcmp(argv[2], "arp")  == 0) filter = ETHERTYPE_ARP;
  else if (strcmp(argv[2], "ipv6") == 0) filter = ETHERTYPE_IPV6;
  else {
    fprintf(stderr, "Protocollo non supportato: %s\n", argv[2]);
    close(sock);
    exit(EXIT_FAILURE);
  }
}
\end{lstlisting}

\begin{lstlisting}[caption={Applicazione filtro nel loop}]
uint16_t proto = ntohs(eth->h_proto);
if (filter != 0 && proto != filter) continue;
\end{lstlisting}

% ---------------------------------------------------------------------------
\section{Domande di approfondimento}
% ---------------------------------------------------------------------------

\subsection*{Q1. Differenza tra \texttt{SOCK\_RAW} e \texttt{SOCK\_DGRAM} con \texttt{AF\_PACKET}}

Con \texttt{SOCK\_RAW} il kernel consegna l'intero frame Ethernet così come viaggia sul
cavo, header MAC incluso. È il livello usato dagli sniffer come Wireshark 
che scende a questo livello tramite libpcap.

Con \texttt{SOCK\_DGRAM} il kernel rimuove l'header Ethernet prima della consegna,
fornendo solo il payload. \texttt{SOCK\_DGRAM} è adatto a messaggi brevi senza
connessione (es.\ server DNS) dove ordine e affidabilità non sono garantiti.
Nel contesto di un sniffer non sarebbe utile, senza header Ethernet non si vedrebbero
i MAC address né l'EtherType.

\subsection*{Q2. Comportamento senza \texttt{bind()}}

Senza \texttt{bind()}, il socket riceve frame da tutte le interfacce con traffico attivo,
non solo da quella passata come argomento.

Non sono riuscito ad effettuare la prova pratica poiché l'unica interfaccia
collegata alla rete che ho attualmente è la wlan0 mentre eth0 è DOWN come
mostrato nell'output del comando "ip a" sopra citato.

\end{document}
